\chapter*{Abstract}
\addcontentsline{toc}{chapter}{Abstract}
The exponential proliferation of embedding-based AI applications has created an urgent demand for scalable, high-performance vector databases. Existing solutions typically require expensive server-grade infrastructure, rendering vector search inaccessible for organizations with limited resources. This paper presents Vectron, a distributed vector database designed to democratize vector search through efficient execution on commodity hardware, from older server-class machines to modern enterprise servers.

Vectron achieves this through an architecture combining Raft-based consensus for metadata management (via Dragonboat), Hierarchical Navigable Small World (HNSW) indexing with SIMD-optimized approximate nearest neighbor search, and int8 quantization reducing memory requirements by 75 percent. The system implements AVX2/AVX-512-accelerated distance computations, a two-stage search pipeline with hot and cold index tiering, and a multi-tier caching strategy comprising gateway-level TinyLFU caching and optional distributed Redis/Valkey caching.

The architecture comprises five core microservices communicating via gRPC: an API Gateway handling request routing and caching, a Placement Driver managing cluster metadata through Raft consensus, Worker nodes hosting shard replicas with HNSW indexing, an Authentication Service for secure access control, and an optional Reranker service for result refinement. The implementation supports asynchronous indexing, search-only replicas fed by WAL streaming, worker-level search batching, and configurable cache tiers.

Benchmark testing demonstrates that Vectron achieves high insertion throughput (up to 2,400+ vectors/sec for lower dimensions), low query latency (P50 of 3-5ms for top-10 queries), and robust scalability (260+ queries/sec with 100 concurrent clients) on commodity infrastructure. These results demonstrate that enterprise-grade vector search performance is achievable on older commodity hardware without infrastructure budget constraints.